%==============================================================
%  lecons/ensembles/construction_R.tex — Construction de R
%==============================================================

\lecontitle{Construction de \texorpdfstring{$\mathbb{R}$}{R}}{Théorie des ensembles — Complétion de \texorpdfstring{$\mathbb{Q}$}{Q}}

%=============================================================
%  § 1 — DEUX APPROCHES
%=============================================================
\secnum{navyblue}{1}{Deux constructions classiques}

\rubriq{Le problème}
$\mathbb{Q}$ est dense mais incomplet. On veut le compléter : obtenir un corps
ordonné archimédien \emph{complet}, c'est-à-dire dans lequel toute suite de
Cauchy converge. Il existe deux méthodes standards.

\begin{tcolorbox}[thmbox]
  \textbf{Construction A — Coupures de Dedekind (1872).}\enspace%
  \index{coupures de Dedekind}\index{réels!construction par coupures}
  Une \emph{coupure} est un couple $(L, U)$ de sous-ensembles non vides de $\mathbb{Q}$
  vérifiant :
  \begin{enumerate}
    \item $L \cup U = \mathbb{Q}$ et $L \cap U = \emptyset$.
    \item $L$ est un \emph{segment initial} : $x \in L$ et $y < x \Rightarrow y \in L$.
    \item $L$ n'a pas de maximum.
  \end{enumerate}
  On pose $\mathbb{R} = \{\,L\mid (L,U)\text{ coupure}\,\}$. L'ensemble $U$ est alors
  entièrement déterminé par $L = \mathbb{Q}\setminus U$ : dans toute la suite, on
  \emph{identifie la coupure à son segment initial $L$}, de sorte que tous les
  énoncés ci-dessous ($\subseteq$, $\sup = \bigcup$, $\hat q = \{r < q\}$) portent
  sur des ensembles de rationnels.
\end{tcolorbox}

\begin{tcolorbox}[thmbox]
  \textbf{Construction B — Suites de Cauchy (Cantor, 1872).}\enspace%
  \index{réels!construction par suites de Cauchy}
  Soit $\mathcal{C}$ l'ensemble des suites de Cauchy dans $\mathbb{Q}$.
  On pose $(a_n) \sim (b_n)$ si $a_n - b_n \to 0$. Alors
  \[
    \mathbb{R} = \mathcal{C}/{\sim}
  \]
  muni des lois terme à terme est un corps complet.
\end{tcolorbox}

\decsep{}

%=============================================================
%  § 2 — CONSTRUCTION PAR COUPURES
%=============================================================
\secnum{navyblue}{2}{Détail de la construction de Dedekind}

\begin{mdframed}[style=proofbar]
  \textit{Ordre.} On pose $\alpha \leq \beta$ si $\alpha \subseteq \beta$ (en tant que
  sous-ensembles de $\mathbb{Q}$). C'est un ordre total sur $\mathbb{R}$.

  \textit{Borne supérieure.} Si $A \subseteq \mathbb{R}$ est non vide et majoré,
  alors $\sup A = \bigcup_{\alpha \in A} \alpha$ est une coupure et c'est le plus
  petit majorant de $A$. C'est la propriété de la \emph{borne supérieure}.

  \textit{Addition.} $\alpha + \beta = \{a + b \mid a \in \alpha,\; b \in \beta\}$.
  On vérifie que c'est une coupure, et que les axiomes de groupe abélien sont satisfaits.

  \textit{Multiplication.} Plus délicate, car il faut gérer les signes.
  Pour $\alpha, \beta > 0$ (c'est-à-dire $\alpha, \beta \supsetneq \hat 0$), on pose
  \[
    \alpha \cdot \beta = \mathbb{Q}^- \cup \{\, ab \mid a \in \alpha,\ b \in \beta,\ a, b > 0 \,\},
  \]
  soit tous les rationnels $\leq 0$ réunis aux produits de représentants strictement
  positifs des deux coupures (on vérifie que l'ensemble obtenu est une coupure~:
  fermé vers le bas, distinct de $\mathbb{Q}$ car majoré par tout produit de
  majorants, et sans maximum). On étend ensuite à $\alpha, \beta$ quelconques par les règles des signes :
  $\alpha\cdot 0 = 0$, et pour $\alpha < 0$ ou $\beta < 0$ on pose
  $\alpha\cdot\beta = \pm\bigl(|\alpha|\cdot|\beta|\bigr)$ selon le signe attendu,
  avec $|\gamma| = \max(\gamma, -\gamma)$.

  \textit{Plongement de $\mathbb{Q}$.} À $q \in \mathbb{Q}$ on associe la coupure
  $\hat{q} = \{r \in \mathbb{Q} \mid r < q\}$. \hfill$\blacksquare$
\end{mdframed}

\decsep{}

%=============================================================
%  § 3 — CONSTRUCTION PAR SUITES DE CAUCHY
%=============================================================
\secnum{navyblue}{3}{Détail de la construction de Cantor}

\rubriq{Étape 1 — L'anneau $\mathcal{C}$}

\begin{mdframed}[style=proofbar]
  Soit $\mathcal{C} = \{(a_n)_{n\in\mathbb{N}} \in \mathbb{Q}^{\mathbb{N}} \mid (a_n)
  \text{ est de Cauchy}\}$ l'ensemble des suites de Cauchy à valeurs dans $\mathbb{Q}$.
  On vérifie que $\mathcal{C}$ est un sous-anneau de $\mathbb{Q}^{\mathbb{N}}$ :

  \begin{itemize}
    \item Si $(a_n)$ et $(b_n)$ sont de Cauchy, alors $(a_n + b_n)$ l'est
    (inégalité triangulaire : $|(a_n+b_n)-(a_m+b_m)| \leq |a_n-a_m|+|b_n-b_m|$).
    \item Si $(a_n)$ et $(b_n)$ sont de Cauchy, elles sont bornées
    ($|a_n| \leq M$, $|b_n| \leq N$). Alors
    $|a_nb_n - a_mb_m| \leq M|b_n-b_m| + N|a_n-a_m| \to 0$,
    donc $(a_nb_n)$ est de Cauchy.
  \end{itemize}
  \hfill$\blacksquare$
\end{mdframed}

\rubriq{Étape 2 — La relation d'équivalence}

\begin{mdframed}[style=proofbar]
  On définit $(a_n) \sim (b_n)$ si $a_n - b_n \to 0$ dans $\mathbb{Q}$.
  C'est bien une relation d'équivalence (réflexivité et symétrie sont immédiates ;
  transitivité : $a_n - c_n = (a_n - b_n) + (b_n - c_n) \to 0$).

  \textit{Compatibilité avec les lois.}
  Si $(a_n)\sim(a_n')$ et $(b_n)\sim(b_n')$, alors
  $(a_n+b_n)-(a_n'+b_n') = (a_n-a_n')+(b_n-b_n') \to 0$ ($+$ bien fondé).
  Pour $\times$ : $(a_nb_n - a_n'b_n') = a_n(b_n-b_n') + b_n'(a_n-a_n') \to 0$
  car $(a_n)$ et $(b_n')$ sont bornées (étant de Cauchy) tandis que
  $(b_n-b_n'),(a_n-a_n')\to 0$ ($\times$ bien fondé).
  \hfill$\blacksquare$
\end{mdframed}

\rubriq{Étape 3 — $\mathbb{R} = \mathcal{C}/{\sim}$ est un corps}

\begin{mdframed}[style=proofbar]
  On note $[(a_n)]$ la classe d'équivalence de $(a_n)$, avec lois
  $[(a_n)]+[(b_n)] = [(a_n+b_n)]$ et $[(a_n)]\cdot[(b_n)]=[(a_nb_n)]$.

  \textit{Neutre additif :} $[(0)]$ (suite constante nulle).
  \textit{Opposé de $[(a_n)]$ :} $[(-a_n)]$.

  \textit{Inverse multiplicatif.}
  Soit $[(a_n)] \neq [(0)]$, i.e.\ $(a_n) \not\to 0$.
  Alors il existe $\varepsilon > 0$ et $N \in \mathbb{N}$ tels que
  $|a_n| > \varepsilon$ pour tout $n \geq N$.
  (En effet, $(a_n)\not\to 0$ fournit un $\varepsilon>0$ et une infinité d'indices
  $m$ tels que $|a_m| > 2\varepsilon$ ; comme $(a_n)$ est de Cauchy, $|a_n-a_m|<\varepsilon$
  pour $n,m\geq N$, d'où $|a_n| \geq |a_m| - |a_n-a_m| > \varepsilon$ dès $n\geq N$.)
  On pose
  \[
    b_n =
    \begin{cases}
      1 & \text{si } n < N, \\
      1/a_n & \text{si } n \geq N.
    \end{cases}
  \]
  Pour $n, m \geq N$ : $|b_n - b_m| = \dfrac{|a_n - a_m|}{|a_n||a_m|}
  \leq \dfrac{|a_n-a_m|}{\varepsilon^2} \to 0$,
  donc $(b_n) \in \mathcal{C}$.
  Et $a_nb_n = 1$ pour $n \geq N$, donc $[(a_n)]\cdot[(b_n)] = [(1)]$.
  Ainsi tout élément non nul est inversible : $\mathbb{R}$ est un corps.
  \hfill$\blacksquare$
\end{mdframed}

\rubriq{Étape 4 — Ordre et plongement de $\mathbb{Q}$}

\begin{mdframed}[style=proofbar]
  \textit{Ordre.} On dit que $[(a_n)] > 0$ s'il existe $\varepsilon > 0$
  et $N$ tels que $a_n > \varepsilon$ pour tout $n \geq N$.
  (Bien fondé : si $(a_n)\sim(a_n')$ et $a_n > \varepsilon$ pour $n\geq N$,
  alors $|a_n-a_n'| < \varepsilon/2$ pour $n$ assez grand, donc $a_n' > \varepsilon/2$.)
  On pose $\alpha \leq \beta$ si $\beta - \alpha \geq 0$.
  C'est un ordre total compatible avec les lois.

  \textit{Plongement.} $\iota : \mathbb{Q} \to \mathbb{R},\; q \mapsto [(q,q,q,\ldots)]$
  est un homomorphisme de corps ordonné injectif.
  On identifie $\mathbb{Q}$ à son image.
  \hfill$\blacksquare$
\end{mdframed}

\rubriq{Étape 5 — $\mathbb{R}$ est complet}

\begin{mdframed}[style=proofbar]
  \textbf{Théorème.}\enspace Toute suite de Cauchy dans $\mathbb{R}$ converge.

  \medskip
  Soit $(\alpha^{(k)})_{k\geq 1}$ une suite de Cauchy dans $\mathbb{R}$,
  avec $\alpha^{(k)} = [(a_n^{(k)})_{n\geq 0}]$.

  \textit{Extraction d'une suite diagonale.}
  Pour chaque $k$, $(a_n^{(k)})$ est de Cauchy dans $\mathbb{Q}$, donc
  il existe $N_k \in \mathbb{N}$ tel que
  \[
    n, m \geq N_k \;\Rightarrow\; |a_n^{(k)} - a_m^{(k)}| < \tfrac{1}{k}.
  \]
  On pose $c_k = a_{N_k}^{(k)} \in \mathbb{Q}$.

  \textit{$(c_k)$ est de Cauchy dans $\mathbb{Q}$.}
  Pour $k, \ell$ grands, $|\alpha^{(k)} - \alpha^{(\ell)}| < \varepsilon/3$
  (hypothèse Cauchy sur $(\alpha^{(k)})$), et $|\alpha^{(k)} - \iota(c_k)| \leq 1/k$
  (par construction de $c_k$). L'inégalité triangulaire donne $|c_k - c_\ell| < \varepsilon$
  pour $k, \ell$ assez grands.

  \textit{Convergence.} Posons $\lambda = [(c_k)] \in \mathbb{R}$.
  Alors $|\alpha^{(k)} - \lambda| \leq |\alpha^{(k)} - \iota(c_k)| + |\iota(c_k) - \lambda| \to 0$,
  car $|\alpha^{(k)} - \iota(c_k)| \leq 1/k \to 0$ et $|\iota(c_k) - \lambda|\to 0$
  puisque $\lambda$ est précisément la classe de la suite de Cauchy $(c_k)$.
  \hfill$\blacksquare$
\end{mdframed}

\decsep{}

%=============================================================
%  § 4 — PROPRIÉTÉS FONDAMENTALES
%=============================================================
\secnum{navyblue}{4}{Caractérisation axiomatique de $\mathbb{R}$}

\begin{tcolorbox}[thmbox]
  \textbf{Théorème (unicité).}\enspace%
  \index{réels!axiomatique}
  À isomorphisme de corps ordonnés près, il existe un \emph{unique} corps ordonné
  complet archimédien. C'est $\mathbb{R}$.

  \medskip
  \textbf{Propriétés équivalentes à la complétude.}
  Les énoncés suivants sont tous équivalents pour un corps ordonné archimédien :
  \begin{enumerate}
    \item Toute suite de Cauchy converge.
    \item Tout sous-ensemble non vide majoré admet une borne supérieure.
    \item Toute suite de segments emboîtés $[a_n, b_n] \supseteq [a_{n+1}, b_{n+1}]$
          avec $b_n-a_n\to 0$ admet un point commun (segments emboîtés).
    \item Toute suite croissante majorée converge (théorème de la limite monotone).
  \end{enumerate}
\end{tcolorbox}

\rubriq{Non dénombrabilité}

\begin{mdframed}[style=proofbar]
  \textbf{Théorème de Cantor.}\enspace%
  \index{non dénombrabilité de $\mathbb{R}$}
  $\mathbb{R}$ n'est pas dénombrable : $|\mathbb{R}| = 2^{\aleph_0} =: \mathfrak{c}$
  (puissance du continu).

  \textit{Preuve (argument diagonal).}
  Il suffit de montrer que le segment $[0,1]$ n'est pas dénombrable, ce qui
  entraîne \emph{a fortiori} la non-dénombrabilité de $\mathbb{R}$. Supposons donc
  $[0,1]$ dénombrable : $[0,1] = \{x_1, x_2, \ldots\}$.
  Construisons $y \in [0,1]$ différent de chaque $x_n$ : si le $n$-ième chiffre
  décimal de $x_n$ est $d$, on pose le $n$-ième chiffre décimal de $y$ égal à
  $5$ si $d \neq 5$, et $6$ sinon. En n'utilisant que les chiffres $5$ et $6$,
  $y$ admet un développement décimal unique (pas de phénomène $0{,}999\ldots=1$),
  et $y \in [0,1]$ mais $y \neq x_n$ pour tout $n$ (ils diffèrent au rang $n$) :
  contradiction. Cet argument n'établit que la non-dénombrabilité ; l'égalité
  $|\mathbb{R}| = 2^{\aleph_0}$ s'obtient ensuite en comparant $\mathbb{R}$ et
  $\{0,1\}^{\mathbb{N}}$ via les développements binaires, puis en appliquant le
  théorème de Cantor--Bernstein.
  \hfill$\blacksquare$
\end{mdframed}

\decsep{}

%=============================================================
%  § 5 — EXEMPLES, PIÈGES ET EXERCICES
%=============================================================
\secnum{exgreen}{5}{Exemples, pièges et exercices}

\rubriq{Exemples}

\begin{mdframed}[style=exbar]
  \textbf{1. Nombres algébriques et transcendants.}\enspace%
  \index{nombre transcendant}
  Les réels algébriques sur $\mathbb{Q}$ forment un sous-corps dénombrable de $\mathbb{R}$.
  Donc «\,la plupart\,» des réels sont transcendants ($e$ par Hermite 1873,
  $\pi$ par Lindemann 1882).

  \smallskip\noindent
  \textbf{2. Mesure de Lebesgue.}\enspace%
  $\mathbb{Q}$ est de mesure nulle dans $\mathbb{R}$ : $\lambda(\mathbb{Q}) = 0$.
  Cela illustre que densité et mesure positive sont des notions indépendantes.

  \smallskip\noindent
  \textbf{3. Topologie.}\enspace%
  $\mathbb{R}$ est localement compact, $\sigma$-compact, connexe, et localement
  connexe. C'est, à homéomorphisme près, l'unique variété topologique connexe
  de dimension $1$, sans bord, non compacte et à base dénombrable.
\end{mdframed}

\vspace{6pt}
\rubriq{Pièges}

\begin{mdframed}[style=cexbar]
  \textbf{Les deux constructions donnent le même objet.}\enspace%
  Les constructions de Dedekind et de Cantor-Cauchy produisent des corps
  ordonnés complets archimédiens, donc isomorphes par unicité.
  La «\,meilleure\,» construction dépend du contexte.

  \smallskip\noindent
  \textbf{Hypothèse du continu.}\enspace%
  On a $\aleph_0 < \mathfrak{c}$, mais la question de savoir s'il existe
  un cardinal strictement entre $\aleph_0$ et $\mathfrak{c}$ est
  \emph{indécidable} dans ZFC (Cohen, 1963).
\end{mdframed}

\vspace{6pt}
\exolabel

\begin{mdframed}[style=exobar]
  \begin{enumerate}
    \item Vérifier que la borne supérieure de coupures $\sup A = \bigcup_{\alpha\in A}\alpha$
          est bien une coupure de Dedekind.

    \item Montrer que les quatre propriétés de complétude (Cauchy, borne sup,
          intervalles emboîtés, suite monotone majorée) sont équivalentes dans
          un corps ordonné archimédien.

    \item Montrer que $\mathbb{R}\setminus\mathbb{Q}$ est non dénombrable.

    \item Montrer que la complétion de $(\mathbb{Q}, |\cdot|_p)$ pour la valeur
          absolue $p$-adique est $\mathbb{Q}_p$, qui n'est \emph{pas} archimédien.

    \item (Théorème de Baire) Montrer que $\mathbb{R}$ n'est pas réunion
          dénombrable d'ensembles fermés d'intérieur vide.
  \end{enumerate}
\end{mdframed}

\leconfooter{R.\ Dedekind (1872) — G.\ Cantor (1872) — unicité due à E.\ Huntington (1903)}
